\documentclass[11pt]{article}
\usepackage[utf8]{inputenc}
\usepackage{amsmath,amssymb,amsthm}
\usepackage[margin=1in]{geometry}
\usepackage{hyperref}
\usepackage{xcolor}
\usepackage{newunicodechar}
\newunicodechar{ℝ}{\ensuremath{\mathbb{R}}}
\newunicodechar{⟨}{\ensuremath{\langle}}
\newunicodechar{⟩}{\ensuremath{\rangle}}
\newunicodechar{Σ}{\ensuremath{\Sigma}}
\newunicodechar{∇}{\ensuremath{\nabla}}
\newunicodechar{φ}{\ensuremath{\varphi}}
\newunicodechar{ψ}{\ensuremath{\psi}}
\newunicodechar{·}{\ensuremath{\cdot}}
\newunicodechar{≠}{\ensuremath{\neq}}
\title{Same-coordinate covariance under a separable Gibbs measure}
\author{Tide retrospective, laplace seabed}
\date{2026-09-06}
\begin{document}
\sloppy
\maketitle
\section{Setting}
The laplace seabed's separable-potential module already proves that, under a product Gibbs measure
$L(x,y) = U(x)+V(y)$, the covariance between a function of the first coordinate and a function of the
second vanishes. This tide adds the diagonal companion: the covariance of two observables of the
\emph{same} coordinate is inherited unchanged from the corresponding one-dimensional marginal. The two
facts together say the covariance matrix is block-diagonal for separable potentials --- exactly the
structural input the project's Stage-3 multivariate law $\mathrm{Cov}_t[\varphi,\psi] = (1/t)\langle
\nabla\varphi, \Sigma\nabla\psi\rangle$ needs when $\Sigma$ is itself separable. It is a deliberate
step off the (now-complete) monomial cumulant ladder toward a different kind of result.
\section{The thing formalised}
For a separable potential and both marginal partition functions nonzero,
\[
  \mathrm{Cov}_{2D}(f\circ\mathrm{fst}, g\circ\mathrm{fst}) = \mathrm{Cov}_U(f,g),\qquad
  \mathrm{Cov}_{2D}(f\circ\mathrm{snd}, g\circ\mathrm{snd}) = \mathrm{Cov}_V(f,g).
\]
The Lean names are \texttt{gibbsCov\_addSeparable\_fst\_fst\_eq} and its second-coordinate twin, in a
new leaf file \texttt{SeparableSameCoordCov.lean}.
\section{Proof strategy}
The covariance unfolds to $\langle fg\rangle - \langle f\rangle\langle g\rangle$. Each Gibbs
expectation of a same-coordinate observable is computed by the existing separable-observable
factorisation, pairing the observable with the constant $1$ in the free coordinate: $\langle
(fg)(z_1)\rangle_{2D} = \langle fg\rangle_U\cdot\langle 1\rangle_V = \langle fg\rangle_U$, since the
Gibbs expectation of the constant $1$ is $1$. Three such collapses (for $fg$, $f$, $g$) rewrite the
2D covariance into the 1D one; the identity then closes definitionally. This is the same collapse
idiom the cross-covariance-vanishing proof already used, run three times on one coordinate.
\section{Roadblocks and resolutions}
The mathematics compiled on the first attempt. The one process lesson was about \emph{where} to put
the theorems. I first appended them to the core \texttt{AddSeparable.lean}; the proofs checked, but
editing a widely-imported file forces a rebuild of every dependent, and those dependents surfaced
their own pre-existing missing-header linter warnings --- a cascade of noise unrelated to this tide.
Relocating the additions to a new leaf file that imports \texttt{AddSeparable} leaves the core file
untouched, so nothing downstream rebuilds and the tide's build is warning-clean.
\section{What was Mathlib, what was new}
No new Mathlib input --- the entire proof rests on the seabed's own separable-observable factorisation
and the constant-observable expectation. New here is the same-coordinate inheritance, which completes
the block-diagonal covariance picture for separable potentials. No GPT consult was needed.
\end{document}
